\chapter{Encerramento cientifico}

\section{Narrativa final}

O fechamento da pesquisa deve ser apresentado como uma investigacao das
condicoes em que compressao numerica e estrutural em Transformers se convertem
em beneficios fisicos reais de memoria e desempenho. A sequencia experimental
parte de equivalencia conceitual, passa por alteracao numerica, compressao
fisica, kernels fisicos, microbenchmarks, modelos OPT, qualidade, escala,
sensibilidade e configuracoes hibridas.

O resultado central e que reduzir bits ou pesos nao implica automaticamente
acelerar inferencia. Representacao fisica e kernel sao fundamentais; overheads
podem eliminar ganhos; resultados locais em microbenchmarks nao necessariamente
predizem comportamento ponta a ponta; e speedup sem qualidade aceitavel nao e
configuracao util.

\section{Hipoteses finais}

\begin{longtable}{@{}p{0.08\linewidth}p{0.22\linewidth}p{0.60\linewidth}@{}}
\toprule
\textbf{Hip.} & \textbf{Classificacao} & \textbf{Resumo da evidencia} \\
\midrule
H1 & sustentada & Reducao de precisao ou pesos nao implicou automaticamente reducao de latencia. \\
H2 & sustentada & Compressao numerica, compressao fisica e aceleracao fisica apareceram como fenomenos distintos. \\
H3 & sustentada & Aceleracao dependeu de representacao, kernel e backend. \\
H4 & nao sustentada & O crossover fisico de operacoes isoladas nao encontrou ganho confirmado. \\
H5 & parcialmente sustentada & Workload e escala influenciaram os resultados, mas nao garantiram ganho util. \\
H6 & sustentada & Pruning denso revelou redundancia numerica sem ganho fisico automatico. \\
H7 & parcialmente sustentada & 2:4 acionou sparse fisico, mas magnitude pruning degradou qualidade. \\
H8 & sustentada & Weight-only reduziu memoria/armazenamento e preservou qualidade sem acelerar. \\
H9 & sustentada & Componentes e layers apresentaram sensibilidades diferentes a quantizacao. \\
H10 & nao sustentada & A politica de sensibilidade dos modelos menores nao transferiu ao OPT-6.7B. \\
H11 & nao sustentada & Nenhum hibrido combinou qualidade preservada e speedup confirmado no maior modelo. \\
H12 & parcialmente sustentada & Distribuicao e outliers influenciaram comportamento sem mudar a conclusao central. \\
\bottomrule
\end{longtable}

\section{Resultados hibridos}

Nas configuracoes avaliadas, a seletividade por componente foi capaz de
preservar qualidade no OPT-1.3B, mas esse comportamento nao se transferiu para
o OPT-6.7B. Portanto, a sensibilidade observada em modelos menores nao se
mostrou diretamente transferivel entre escalas.

No OPT-1.3B, os finalistas hibridos com MLP quantizado e atencao preservada
apresentaram medianas de desempenho positivas, como 1,110x, 1,131x e 1,108x em
prefill. Contudo, os intervalos de 95\% cruzaram 1,0x, de modo que nao existe
speedup estatisticamente confirmado. No OPT-6.7B, as configuracoes hibridas
avaliadas falharam qualidade, com aumento de perplexidade aproximadamente entre
+157\% e +207\%. Por isso, os benchmarks fisicos dessas candidatas foram
corretamente evitados.

\section{Resultados negativos}

Os resultados negativos sao parte da contribuicao. A ausencia de crossover
confirmado mostra que kernels fisicos podem nao compensar overheads em
operacoes isoladas. O weight-only mostrou compressao e preservacao de qualidade
sem aceleracao no backend avaliado. O 2:4 confirmou sparse fisico, mas a
selecao por magnitude degradou qualidade. O INT8 dinamico acelerou caminhos do
OPT-6.7B, mas com perplexidade fora dos limites. Os hibridos preservaram
qualidade no OPT-1.3B sem speedup confirmado e falharam qualidade no OPT-6.7B.

\section{Lacuna opcional: 2:4 com selecao melhor}

A unica extensao experimental ainda potencialmente relevante e melhorar a
selecao de pesos do 2:4. Wanda e a abordagem mais adequada se houver tempo:
e post-training, usa informacao de ativacao, dispensa fine-tuning pesado e pode
preservar o mesmo backend sparse caso a saida seja convertida para 2:4.
SparseGPT tambem e relevante, mas tem custo computacional e complexidade de
integracao maiores. A classificacao final dessa extensao e util mas opcional,
nao essencial para encerrar a iniciacao cientifica.

\section{Decisao de encerramento}

Os dados atuais sao suficientes para encerrar a IC. A pesquisa responde a
pergunta principal com evidencia rastreavel: alterar numeros, comprimir
representacao e acelerar fisicamente sao niveis diferentes de evidencia. O
beneficio final depende de representacao, kernel, workload, qualidade e escala.
